% =====================================================================
%  SCALAR -- Supplementary Material
%  WACV 2027 Submission #2556, Algorithms Track
%
%  STANDALONE: compiles on its own with pdflatex.
%  To match the main paper, replace the \documentclass line and preamble
%  with the wacv class:   \documentclass[10pt,twocolumn,letterpaper]{article}
%                         \usepackage{wacv}          % or \usepackage[review]{wacv}
%  and delete the \geometry line below.
% =====================================================================
\documentclass[10pt,twocolumn,letterpaper]{article}

\usepackage[margin=0.75in]{geometry}
\usepackage{times}
\usepackage{amsmath,amssymb,amsfonts}
\usepackage{graphicx}
\usepackage{booktabs}
\usepackage[table]{xcolor}
\usepackage{multirow}
\usepackage{algorithm}
\usepackage{algpseudocode}
\usepackage{url}
\usepackage[colorlinks,linkcolor=blue,citecolor=blue]{hyperref}

\usepackage{pifont}

\newcommand{\ours}{SCALAR}
\newcommand{\cmark}{\ding{51}}
\newcommand{\xmark}{\ding{55}}

\renewcommand{\thesection}{\Alph{section}}
\renewcommand{\thetable}{S\arabic{table}}
\renewcommand{\thefigure}{S\arabic{figure}}
\renewcommand{\theequation}{S\arabic{equation}}

\title{\Large Query-Conditioned Spherical Centroid Aggregation for\\
       Multimodal Retrieval\\[2mm]
       \large Supplementary Material}
\author{Anonymous WACV Algorithms Track submission\\Paper ID 2556}
\date{}

\begin{document}
\maketitle

\noindent
This supplement is organised as follows. Section~\ref{sec:llm} states our use of
large language models and Section~\ref{sec:repro} our reproducibility protocol.
Section~\ref{sec:algo} gives the aggregation and training procedures in
pseudocode, Section~\ref{sec:impl} the complete hyperparameter set, and
Section~\ref{sec:data} the dataset statistics. Section~\ref{sec:extra} reports
the additional results referenced from the main paper: the components of the
aggregation gain, the full two-stage results under masking, the modality-subset
ladder, a latent-space analysis, a negative result on VGGSound, an
evaluation-environment audit, and the statistical tests behind every count in
the main text. Section~\ref{sec:limits} states the limitations.

% =====================================================================
\section{Use of Large Language Models}
\label{sec:llm}
% =====================================================================

Large language models were used in the preparation of this paper for two
purposes: (i) polishing the writing---grammar, phrasing and consistency of
terminology---in text whose technical content was written by the authors, and
(ii) compressing existing passages so that the paper fits the page limit. They
were not used to generate research ideas, to design or run experiments, to
produce or interpret results, or to write code that produced any reported
number. Every claim, table, figure and analysis in this paper was derived from
experiments the authors designed and executed, and was checked by the authors
against the underlying logs. The authors take full responsibility for the
entire content of the paper, including any text that passed through such a
tool.

% =====================================================================
\section{Reproducibility}
\label{sec:repro}
% =====================================================================

\paragraph{What is released.}
We will release the full training and evaluation code, every configuration file
used to produce a reported number, the scripts that extract each table and
figure directly from the evaluation logs, and the adapter weights for the
reported configuration. Every table in the paper is emitted by a generator that
reads evaluation logs; no number is transcribed by hand, and a cell that cannot
be read from a log is printed as \texttt{MISSING} and aborts generation rather
than being filled in.

\paragraph{Single reported configuration.}
All \ours{} numbers in the paper come from one configuration, applied unchanged
to every benchmark. We do not tune per benchmark, and no benchmark-specific
learning rate, temperature or schedule is used anywhere. The three seeds
(50/51/52) differ only in the seed value. The full-fine-tuning control is that
same configuration with the adapters disabled; the uniform-weights control is
that same configuration with query conditioning disabled. Each ablation in
Table~5 of the main paper removes exactly one component and retunes nothing.

\paragraph{Baseline artifacts.}
GRAM and PMRL are their authors' released checkpoints, evaluated by us.
HyperGRAM releases no checkpoint, so it is trained from its authors' unmodified
code at their published recipe. No baseline number in the measured blocks is
taken from a paper. All methods start from the same VAST foundation checkpoint
and are evaluated in one environment under one protocol.

\paragraph{Determinism.}
The missing-modality masks are fully deterministic and reproducible without
storing them: for clip identifier $c$ and seed $\sigma$ we take
$h = \mathrm{md5}(\sigma \,\|\, c)$, drop a modality iff
$h_{[0:4]}/2^{32} < r$, and select the dropped index as
$h_{[4:8]} \bmod |M|$. Because the dropped index does not depend on $r$, the
masked sets are \emph{nested} across rates: every clip masked at $r$ is masked
at every $r' > r$. All methods receive byte-identical masks, so the comparison
is paired rather than two independent draws.

\paragraph{Evaluation noise.}
Repeated evaluation of identical checkpoints across 15 cells bounds the
observed spread at $0.2$ R@1. We treat this as a floor: no margin at or below
$0.2$ is reported as a result, and every headline margin in the paper is quoted
against it. Training itself is not bit-deterministic (mixed precision, four-way
data-parallel reduction), which is why \ours{} rows are reported as a mean over
three seeds with the standard deviation.

\paragraph{Compute.}
Each \ours{} pretraining run takes approximately 11 hours on four NVIDIA A100
64\,GB GPUs. The complete set of experiments in the paper---three seeds, the
full-fine-tuning control, the uniform-weights control, four objective
ablations, the $5\times5$ missing-modality grid for two methods, and the
baseline evaluations---amounts to roughly 900 GPU-hours.

% =====================================================================
\section{Algorithms}
\label{sec:algo}
% =====================================================================

Algorithm~\ref{alg:score} gives inference. The centroid is query-dependent, but
it is never materialised per query--candidate pair: the within-sample Gram
matrix $[\langle z_m, z_n\rangle]$ is $|M|\times|M|$ and is computed once per
candidate, after which scoring a query against a candidate costs $O(|M|^2)$
regardless of the embedding dimension. A missing modality is simply absent from
$M$; nothing is imputed and no candidate is discarded.

\begin{algorithm}[t]
\caption{\ours{} scoring for one query--candidate pair}
\label{alg:score}
\begin{algorithmic}[1]
\Require query embedding $z^T\!\in\!\mathbb{S}^{d-1}$; observed modality
         embeddings $\{z_m\}_{m\in M}$, $M\subseteq\{V,A,S\}$, $M\neq\emptyset$;
         temperature $\tau_w$
\Ensure  score $s\in[-1,1]$
\State $a_m \gets \langle z^T\!, z_m\rangle$ \Comment{query agreement, $m\in M$}
\State $w_m \gets \dfrac{\exp(a_m/\tau_w)}{\sum_{n\in M}\exp(a_n/\tau_w)}$
       \Comment{softmax over the \emph{observed} set only}
\State $G_{mn} \gets \langle z_m, z_n\rangle$ \Comment{cached per candidate}
\State $s \gets \dfrac{\sum_{m} w_m a_m}
                     {\sqrt{\sum_{m,n} w_m w_n G_{mn}}}$
       \Comment{$=\langle z^T\!, \mu_M(z^T)\rangle$}
\State \Return $s$
\end{algorithmic}
\end{algorithm}

Algorithm~\ref{alg:train} gives one training step. The two views---complete and
reduced---come from a single forward pass: masking is applied after encoding, so
the reduced-arity view costs one extra centroid, not one extra encoder pass.

\begin{algorithm}[t]
\caption{One \ours{} training step}
\label{alg:train}
\begin{algorithmic}[1]
\Require batch of $B$ matched (text, clip) pairs; step $t$; complete modality
         set $\mathcal{K}$ per clip
\State encode all modalities once; L2-normalise
\State $p_{\mathrm{full}} \gets \max(0.5,\, 1 - 0.5\,t/2000)$
\For{each clip in the batch}
  \If{$\mathrm{Bernoulli}(p_{\mathrm{full}})$ \textbf{or} $|\mathcal{K}|<3$}
     \State $\mathcal{M} \gets \mathcal{K}$
  \Else
     \State $m^\dagger \sim \mathrm{Unif}(\mathcal{K})$;\;
            $\mathcal{M} \gets \mathcal{K}\setminus\{m^\dagger\}$
  \EndIf
\EndFor
\State $\mu_{\mathcal{M}}, \mu_{\mathcal{K}} \gets$ query-weighted centroids
       (Alg.~\ref{alg:score}, lines 1--4) for both views
\State $C_{ij} \gets \langle z^T_i, \mu_{\mathcal{M}_j}(z^T_i)\rangle$;\quad
       $\tilde{C}_{ij} \gets \langle z^T_i, \mu_{\mathcal{M}_j}(z^T_j)\rangle$
\State $\mathcal{L}_{\mathrm{align}} \gets
        \tfrac12[\mathrm{InfoNCE}(C/\tau) + \mathrm{InfoNCE}(C^\top\!/\tau)]$
\State $\mathcal{L}_{\mathrm{mask}} \gets
        1-\langle\mu_{\mathcal{M}},\mu_{\mathcal{K}}\rangle
        + (s_{\mathcal{M}}-\mathrm{sg}[s_{\mathcal{K}}])^2$
\If{$t > 500$} \Comment{warm-up gate}
  \State $\mathcal{L}_{\mathrm{sem}}, \mathcal{L}_{\mathrm{unif}} \gets$
         Eqs.~(8), (9) of the main paper, using $\tilde{C}$ for
         $\mathcal{L}_{\mathrm{sem}}$
\Else
  \State $\mathcal{L}_{\mathrm{sem}}, \mathcal{L}_{\mathrm{unif}} \gets 0$
\EndIf
\State $\mathcal{L}_{\mathrm{DAM}} \gets$ hard-negative matching loss, negatives
       drawn from $C$
\State \Return $\mathcal{L}_{\mathrm{align}} + \mathcal{L}_{\mathrm{sem}}
        + \mathcal{L}_{\mathrm{mask}} + 0.1\,\mathcal{L}_{\mathrm{unif}}
        + 0.1\,\mathcal{L}_{\mathrm{DAM}}$
\end{algorithmic}
\end{algorithm}

Two details in Algorithm~\ref{alg:train} are load-bearing and easy to get wrong.
First, $C$ and $\tilde{C}$ differ in \emph{which} text conditions the weights:
in $C$ every candidate is re-aggregated with respect to the query that scores
it, which is what inference does, whereas $\tilde{C}$ aggregates each candidate
once using its own paired text. Using $\tilde{C}$ for retrieval would let the
model win by sharpening weights rather than by learning features, because the
positive pair alone would receive a centroid built for its own text. Second,
only $s_{\mathcal{K}}$ is stop-gradiented in
$\mathcal{L}_{\mathrm{mask}}$---the complete-set score is the fixed reference
that the reduced-arity score is pulled towards, not a symmetric penalty.

% =====================================================================
\section{Implementation Details}
\label{sec:impl}
% =====================================================================

Table~\ref{tab:hparams} lists every hyperparameter of the reported
configuration. All values are shared across the five benchmarks.

\begin{table}[t]
\centering
\caption{Complete hyperparameters for the reported \ours{} configuration.
Identical for every benchmark; the three seeds differ only in the seed value.}
\label{tab:hparams}
\small
\setlength{\tabcolsep}{4pt}
\begin{tabular}{ll}
\toprule
\multicolumn{2}{l}{\emph{Backbone (frozen)}} \\
Initialisation & VAST foundation checkpoint \\
Vision encoder & EVA-CLIP ViT-g/14 \\
Audio encoder & BEATs \\
Text encoder & BERT-base (subtitle + query) \\
Projection dimension $d$ & 512 \\
\midrule
\multicolumn{2}{l}{\emph{Adaptation}} \\
LoRA rank $r$ & 8 \\
LoRA $\alpha$ & 16 \\
LoRA dropout & 0.0 \\
LoRA target modules & query, value projections \\
Trainable parameters & 4{,}763{,}650 \\
\midrule
\multicolumn{2}{l}{\emph{Optimisation}} \\
Optimiser & AdamW, $\beta=(0.9, 0.98)$ \\
Learning rate & $2\times10^{-5}$ \\
Weight decay & 0.01 \\
Gradient-norm clip & 2.0 \\
Schedule & linear, warm-up ratio 0.1 \\
Batch size & 128 \\
Epochs / steps & 5 / 5{,}330 \\
Frames sampled (training) & 2 \\
Precision & fp16, DDP over 4 GPUs \\
Seeds & 50, 51, 52 \\
\midrule
\multicolumn{2}{l}{\emph{Aggregation and losses}} \\
Query-weight temperature $\tau_w$ & 0.1 \\
Contrastive temperature $\tau$ & 0.07 (fixed, not learned) \\
Label smoothing & 0.1 \\
$\mathcal{L}_{\mathrm{sem}}$ weight $\alpha$ & 1.0 \\
$\mathcal{L}_{\mathrm{mask}}$ weight $\beta$ & 1.0 \\
$\mathcal{L}_{\mathrm{unif}}$ weight $\lambda$ & 0.1 \\
$\mathcal{L}_{\mathrm{unif}}$ scale $t$ & 2.0 \\
$\mathcal{L}_{\mathrm{DAM}}$ weight & 0.1 \\
Semantic temperature $\tau^*$ & 0.5 \\
Semantic neighbours kept & top 64 \\
Calibration weight & 1.0 (observed pairs only) \\
Loss warm-up ($\mathcal{L}_{\mathrm{sem}},\mathcal{L}_{\mathrm{unif}}$) & 500 steps \\
\midrule
\multicolumn{2}{l}{\emph{Reduced-arity views}} \\
$p_{\mathrm{full}}$ schedule & $1.0 \rightarrow 0.5$, linear \\
Anneal length & 2{,}000 steps \\
Modalities dropped & exactly 1, uniform over $\mathcal{K}$ \\
Minimum retained & 2 \\
\midrule
\multicolumn{2}{l}{\emph{Inference}} \\
Frames sampled (evaluation) & 8 \\
Rerank depth & top 50 \\
Eval batch size & 64 \\
Test-time mask rates $r$ & 0, 25, 50, 75, 90\% \\
Mask seed & 0 \\
\bottomrule
\end{tabular}
\end{table}

\paragraph{Frame budget.}
We sample 2 frames per clip during continued pretraining and 8 at evaluation.
The compared methods use 8 for both~\cite{cicchetti2025gramian}, so our
adaptation sees a quarter of their visual input per clip while every reported
comparison is made at the same 8 evaluation frames for all methods. The
reduced training budget is applied identically to the reported configuration,
the full-fine-tuning control, the uniform-weights control and every ablation,
so the comparisons among them are unaffected. We swept the training frame
count rather than assuming it was harmless; Section~\ref{sec:frames} reports
the result.

\paragraph{Temperature fixed by construction.}
$\tau$ must not be learnable in this objective. The calibration term of
$\mathcal{L}_{\mathrm{sem}}$ pins the absolute cosine scale to the target
affinity, whereas softmax is shift- and scale-invariant, so a learnable $\tau$
would silently re-absorb the calibration and void the claim. Our configuration
layer raises an error if calibration is requested together with a learnable
temperature.

\paragraph{Aggregation and the semantic objective use different score
matrices.}
$\mathcal{L}_{\mathrm{sem}}$ operates on $\tilde{C}$, not $C$, so the graded
targets shape the embedding geometry independently of the query-conditioned
weighting. This keeps the semantic supervision from interacting with the
mechanism the paper is testing.

% =====================================================================
\section{Data Statistics}
\label{sec:data}
% =====================================================================

\begin{table}[t]
\centering
\caption{Benchmark statistics. Modalities and splits follow the released
baselines. VATEX is evaluated on the 431 test clips that remain downloadable;
absolute VATEX recalls are therefore comparable within our tables but not to
numbers reported on the full test set.}
\label{tab:datastats}
\small
\setlength{\tabcolsep}{4pt}
\begin{tabular}{lccccc}
\toprule
Benchmark & Modalities & Train & Val & Test & Frames \\
\midrule
MSR-VTT & V, A, S & 9{,}000 & -- & 1{,}000 & 8 \\
DiDeMo & V, A & 8{,}394 & 1{,}065 & 1{,}003 & 8 \\
ActivityNet & V, A & 10{,}009 & -- & 4{,}917 & 8 \\
VATEX & V, A, S & 14{,}060 & -- & 431 & 8 \\
AudioCaps & V, A & -- & -- & 700 & 8 \\
\midrule
VGGSound-5K & V, A & -- & -- & 5{,}000 & 8 \\
\bottomrule
\end{tabular}
\end{table}

\paragraph{Training corpus.}
All compared methods continue-pretrain from the same VAST checkpoint on the
same 150k-clip subset of VAST-27M. Of the nominal 150k clips, 136{,}674 remain
downloadable and form the effective training set for every arm in this paper;
five epochs at batch size 128 over that set is 5{,}330 optimizer steps. We
report the effective size because it fixes the training budget: the compared
papers report the nominal 150k figure, and the same source attrition affects
them.

\paragraph{Modality availability.}
The five benchmarks span two arities, which is why arity-invariance is a
practical requirement rather than a formal nicety: MSR-VTT and VATEX are scored
with three gallery modalities and the remaining three benchmarks with two. A
single reported configuration is therefore evaluated at both $|M|=2$ and
$|M|=3$ with no change, and the same construction serves the reduced arities
produced by test-time masking.

% =====================================================================
\section{Additional Results and Analysis}
\label{sec:extra}
% =====================================================================

\subsection{Components of the Aggregation Gain}
\label{sec:gaincomp}

Table~\ref{tab:gain_components} decomposes every cell of the aggregation-gain
table in the main paper into its two terms, so the subtraction is checkable.
Two rows deserve comment.

On ActivityNet, GRAM's best single pathway is $36.8$ against our $34.1$: its
unimodal features are $2.7$ points \emph{better} than ours. Its volume
aggregate is $31.0$ and ours is $34.4$. The volume gives up $5.8$ points
relative to the best input it was handed; the centroid gains $0.3$. AudioCaps
has the same shape---best single pathway $29.0$ for GRAM against $25.7$ for us,
aggregate $22.9$ against $27.1$. This rules out the obvious alternative
explanation for our results: \ours{} is not ahead because low-rank adaptation
produced better encoders. On two of five benchmarks its encoders are measurably
worse. It is ahead because the aggregation step does not discard the
information they provide.

\begin{table*}[t]
\centering
\caption{Components of the aggregation gain: best single-modality pathway R@1,
the method's own aggregate R@1, and their difference, per benchmark.}
\label{tab:gain_components}
\small
\setlength{\tabcolsep}{3pt}
\begin{tabular}{lccccccccccccccc}
\toprule
 & \multicolumn{3}{c}{MSR-VTT} & \multicolumn{3}{c}{DiDeMo} & \multicolumn{3}{c}{ActivityNet} & \multicolumn{3}{c}{VATEX} & \multicolumn{3}{c}{AudioCaps} \\
\cmidrule(lr){2-4}\cmidrule(lr){5-7}\cmidrule(lr){8-10}\cmidrule(lr){11-13}\cmidrule(lr){14-16}
Method & best & agg. & gain & best & agg. & gain & best & agg. & gain & best & agg. & gain & best & agg. & gain \\
\midrule
GRAM$^{\star}$ & 42.1 & 38.7 & -3.4 & 32.0 & 28.2 & -3.8 & 36.8 & 31.0 & -5.8 & 77.5 & 75.6 & -1.9 & 29.0 & 22.9 & -6.1 \\
PMRL$^{\star}$ & 43.7 & 31.5 & -12.2 & 38.3 & 28.7 & -9.6 & 39.8 & 30.0 & -9.8 & 81.2 & 53.6 & -27.6 & 32.0 & 33.4 & +1.4 \\
HyperGRAM$^{\dagger}$ & 42.5 & 39.1 & -3.4 & 32.1 & 32.0 & -0.1 & 36.3 & 34.3 & -2.0 & 78.0 & 78.0 & +0.0 & -- & -- & -- \\
\ours{}, uniform weights & 41.1 & 34.1 & -7.0 & 34.5 & 27.0 & -7.5 & 30.0 & 26.0 & -4.0 & 80.3 & 69.1 & -11.2 & 28.6 & 28.8 & +0.2 \\
\textbf{\ours{}, query-weighted (ours)} & 41.2 & 45.2 & +4.0 & 35.3 & 34.2 & -1.1 & 34.1 & 34.4 & +0.3 & 81.2 & 81.7 & +0.5 & 25.7 & 27.1 & +1.4 \\
\bottomrule
\end{tabular}
\end{table*}

\subsection{Two-Stage Results Under Masking}
\label{sec:twostage}

Table~\ref{tab:missing_twostage} reports the reranked metric on the identical
masked cells of the main paper's Table~3, alongside each method's own video-only
cosine at $r{=}90\%$. It is the evidence for reporting the representation stage
in the main text.

Each method reranks with the cross-encoder in its own checkpoint, all descended
from the same VAST foundation model, and none of them was ever trained on an
incomplete modality set: the reranking branch conditions on the full encoder
outputs, and the baselines' pipelines discard incomplete clips before training.
As $r$ grows, every method's two-stage score converges onto its own video-only
cosine---at $r{=}90\%$, $33.2$ against $32.1$ for \ours{} and $32.9$ against
$32.8$ for GRAM on MSR-VTT---and the between-method spread collapses from $2.3$
points at $r{=}0$ to within evaluation noise. Beyond the $r{=}0$ anchor the
two-stage number therefore measures the reranker's behaviour on zeroed inputs
rather than the aggregation geometries under comparison. We report it here in
full rather than omitting it: on DiDeMo, ActivityNet and AudioCaps the reranked
ordering under masking favours GRAM, and that is a consequence of the same
collapse, not a contradiction of the representation-level result.

\begin{table}[t]
\centering
\caption{Two-stage (reranked) R@1 on the identical masked cells of the main
paper's Table~3, with each method's own video-only cosine at $r{=}90\%$. The
reranking stage converges every method onto that pathway as $r$ grows, erasing
the aggregator differences: end-to-end robustness in a two-stage pipeline is
capped by the reranking stage.}
\label{tab:missing_twostage}
\small
\setlength{\tabcolsep}{4pt}
\begin{tabular}{llcccccc}
\toprule
 & Method & 0\% & 25\% & 50\% & 75\% & 90\% & cos$_{TV}$@90\% \\
\midrule
MSR-VTT & \textbf{\ours{}} & \textbf{54.8} & \textbf{47.8} & \textbf{41.1} & \textbf{37.0} & \textbf{33.2} & 32.1 \\
 & GRAM$^{\star}$ & 52.5 & 47.0 & 40.9 & 35.4 & 32.9 & 32.8 \\
\rowcolor{gray!10} & \emph{margin} & +2.3 & +0.8 & +0.2 & +1.6 & +0.3 &  \\
\midrule
DiDeMo & \textbf{\ours{}} & \textbf{51.3} & 44.2 & 38.9 & 32.1 & 29.5 & 23.2 \\
 & GRAM$^{\star}$ & 50.7 & \textbf{45.0} & \textbf{39.8} & \textbf{32.7} & \textbf{29.8} & 21.1 \\
\rowcolor{gray!10} & \emph{margin} & +0.6 & -0.8 & -0.9 & -0.6 & -0.3 &  \\
\midrule
ActivityNet & \textbf{\ours{}} & 55.8 & 47.9 & 40.0 & 33.8 & 30.7 & 23.7 \\
 & GRAM$^{\star}$ & \textbf{56.3} & \textbf{48.3} & \textbf{41.4} & \textbf{35.5} & \textbf{32.3} & 24.8 \\
\rowcolor{gray!10} & \emph{margin} & -0.5 & -0.4 & -1.4 & -1.7 & -1.6 &  \\
\midrule
VATEX & \textbf{\ours{}} & \textbf{90.5} & \textbf{85.2} & \textbf{73.3} & \textbf{61.9} & 57.3 & 58.5 \\
 & GRAM$^{\star}$ & 90.0 & 82.8 & 72.9 & \textbf{61.9} & \textbf{57.8} & 58.0 \\
\rowcolor{gray!10} & \emph{margin} & +0.5 & +2.4 & +0.4 & +0.0 & -0.5 &  \\
\midrule
AudioCaps & \textbf{\ours{}} & \textbf{35.2} & 24.6 & 19.2 & 13.2 & 9.9 & 10.2 \\
 & GRAM$^{\star}$ & 32.2 & \textbf{26.0} & \textbf{21.0} & \textbf{15.5} & \textbf{11.4} & 9.5 \\
\rowcolor{gray!10} & \emph{margin} & +3.0 & -1.4 & -1.8 & -2.3 & -1.5 &  \\
\bottomrule
\end{tabular}
\end{table}

\subsection{Modality-Subset Ladder}
\label{sec:ladder}

Table~\ref{tab:subset_ladder} evaluates the same checkpoints with a fixed
modality set for the whole gallery---$\{V\}$, $\{V,A\}$, $\{V,A,S\}$---on the
two benchmarks that carry three gallery modalities. It answers a different
question from the missing-modality experiment and the two should not be
conflated. The ladder fixes arity for the entire gallery and therefore measures
each modality's \emph{contribution}; because arity is constant within a run, an
arity-dependent score is a global rescaling that leaves the ranking unchanged.
The masking experiment instead degrades clips \emph{independently}, so
candidates of different arity compete in one ranked list, which is the setting
in which cross-arity comparability becomes a ranking property rather than a
cosmetic one.

The ladder makes the mechanism visible constructively: every modality \ours{}
adds helps or holds, while every modality the volume adds hurts. The $\{V\}$
rung doubles as a control---a one-element aggregate must equal the text--video
cosine for both geometries, which we assert at generation.

\begin{table}[t]
\centering
\caption{Modality-subset ladder: each method's own aggregation score
(T$\rightarrow$V R@1) as modalities are added at test time, using the same
checkpoints as the main tables. $\{V,A,S\}$ \emph{is} the main-table cell.}
\label{tab:subset_ladder}
\small
\setlength{\tabcolsep}{6pt}
\begin{tabular}{llccc c}
\toprule
 & Method & $\{V\}$ & $\{V,A\}$ & $\{V,A,S\}$ & $\Delta_{V \to VAS}$ \\
\midrule
MSR-VTT & \textbf{\ours{}} & 41.2 & 42.1 & 45.2 & +4.0 \\
 & GRAM$^{\star}$ & 42.1 & 39.7 & 38.7 & -3.4 \\
\midrule
VATEX & \textbf{\ours{}} & 81.2 & 81.9 & 81.7 & +0.5 \\
 & GRAM$^{\star}$ & 77.5 & 75.6 & 75.6 & -1.9 \\
\bottomrule
\end{tabular}
\end{table}

\subsection{Latent-Space Analysis}
\label{sec:tsne}

Figure~\ref{fig:tsne} visualises the embedding space of three models that share
an initialisation: the VAST foundation checkpoint both adapted models start
from, the fully fine-tuned GRAM checkpoint, and \ours{}. Because t-SNE is
qualitative and the display is a subset, we separate the two layers explicitly.
The numbers printed on each panel are computed in the \emph{original} embedding
space using cosine distance and over \emph{all eight} dumped classes---never
over the displayed subset and never in t-SNE coordinates. The display shows the
three classes with the largest gain of both adapted models over the shared
initialisation on text-centroid cosine, a deterministic rule fixed before
looking at the panels.

Mean text-centroid cosine---how tightly a class's text embedding anchors the
mean of that class's video and audio embeddings---is $0.34$ for VAST, $0.48$ for
GRAM and $0.44$ for \ours{}. Mean audio silhouette by class is $0.21$, $0.26$
and $0.20$ respectively. The reading is deliberately modest: alignment training
relocates text anchors substantially relative to the shared initialisation, and
4.8M adapter parameters recover most, though not all, of what full fine-tuning
of 1.4B parameters buys on this axis. We report the axis on which we do not
lead rather than selecting one on which we do.

\begin{figure}[t]
\centering
\includegraphics[width=\linewidth]{figures/fig_tsne_latent.pdf}
\caption{Joint t-SNE of text, video and audio embeddings on a VGGSound subset,
for the shared VAST initialisation, the fully fine-tuned GRAM checkpoint and
\ours{}. Stars are class text embeddings, squares video, triangles audio.
Panel annotations are computed in the original cosine space over all eight
classes, not over the displayed three.}
\label{fig:tsne}
\end{figure}

\subsection{VGGSound: A Negative Result and a Protocol Check}
\label{sec:vggsound}

Table~\ref{tab:vggsound} reports audio-anchored classification-as-retrieval on
VGGSound-5K, where \ours{} does not lead. We include it for two reasons.

First, the decomposition locates the deficit. \ours{}'s video pathway is weaker
than the fully fine-tuned baseline's ($31.3$ against $35.0$) while its audio
pathway is slightly stronger ($27.0$ against $26.3$); its fusion still adds
$+3.6$ over its own best pathway, so the aggregation behaves as the rest of the
paper predicts even where the adapted pathways are not competitive. Short
label-style queries---310 class names rather than sentences---are the regime in
which rank-8 adaptation of a frozen backbone is weakest, and we do not claim
otherwise. Notably, the uniform-weights control is \emph{stronger} here
($35.9$), the one setting in the paper where query conditioning does not help;
with a single-word query there is little for the weights to condition on.

Second, the GRAM$^{\star}$ row is the strongest protocol validation in the
paper. The released checkpoint reproduces its published VGGSound number to
within $0.1$ in our environment ($40.5$ measured against $40.6$ published),
which bounds the environment offset on this benchmark far more tightly than the
retrieval benchmarks do.

\begin{table}[t]
\centering
\caption{VGGSound-5K, audio-anchored classification-as-retrieval: 310 class
labels as queries over 5{,}000 clips, Acc@1/Acc@10, no reranking stage---not
comparable to the retrieval tables. \ours{} Acc@1 is three seeds ($\pm$ sd).}
\label{tab:vggsound}
\small
\setlength{\tabcolsep}{6pt}
\begin{tabular}{lcc}
\toprule
Method & Acc@1 & Acc@10 \\
\midrule
GRAM (paper)$^{\S}$ & 40.6 & 78.1 \\
GRAM$^{\star}$ & 40.5 & 77.3 \\
\textbf{\ours{}} & 34.9\tiny{$\pm$0.3} & 71.2 \\
\midrule
\multicolumn{3}{l}{\emph{decomposition} (\ours{} / GRAM$^{\star}$)} \\
\quad video pathway alone & 31.3 / 35.0 & 63.9 / 68.4 \\
\quad audio pathway alone & 27.0 / 26.3 & 63.3 / 66.4 \\
\quad \ours{}, uniform weights & 35.9 & 72.7 \\
\bottomrule
\end{tabular}
\end{table}

\subsection{Training Frame Budget}
\label{sec:frames}

The compared methods sample 8 frames per clip for both training and
inference~\cite{cicchetti2025gramian}; we sample 2 for training and 8 for
inference. Because this is the one recipe element where we depart from the
baselines, we measured its cost directly rather than arguing it away.

Table~\ref{tab:frames} reports two additional pretraining runs of the reported
configuration, changing only the frame budget. At a fixed batch size of 128,
doubling the training frames from 2 to 4 changes the mean two-stage
R@1 by $+0.04$---inside the $0.2$ evaluation-noise floor---and the mean
aggregator R@1 by $-0.22$. Reaching the baselines' 8 frames requires halving
the batch to 64 to fit in memory, and therefore halves the in-batch negatives;
that run trains to completion over 10{,}670 steps and its MSR-VTT validation
tops out at $53.0$ against the reported configuration's $54.8$.

More frames do improve the visual pathway, in the expected direction: the
text--video cosine rises on four of five benchmarks at 4 frames ($41.2
\rightarrow 41.6$ on MSR-VTT, $35.3 \rightarrow 35.5$ on DiDeMo, $34.1
\rightarrow 34.5$ on ActivityNet, $81.2 \rightarrow 81.4$ on VATEX). That
improvement does not survive into the aggregate or the reranked metric. We
therefore report the 2-frame configuration, and note that the frame budget is
not where the gap between aggregation geometries lives.

\begin{table}[t]
\centering
\caption{Training frame budget. All three runs are the reported configuration
with only the frame count (and, for 8 frames, the batch size that makes it fit)
changed. $\bar{\Delta}$ columns are means over the five benchmarks; the
8-frame run has in-training validation only.}
\label{tab:frames}
\small
\setlength{\tabcolsep}{4pt}
\begin{tabular}{lccccc}
\toprule
Training frames & Batch & Steps & MSR-VTT & $\bar{\Delta}$ two-stage & $\bar{\Delta}$ agg. \\
\midrule
2 (reported) & 128 & 5{,}330 & 54.8 & -- & -- \\
4 & 128 & 5{,}330 & 54.2 & +0.04 & -0.22 \\
8 & 64 & 10{,}670 & 53.0 & n/a & n/a \\
\bottomrule
\end{tabular}
\end{table}

\subsection{Evaluation-Environment Audit}
\label{sec:envaudit}

Table~\ref{tab:envaudit} evaluates \emph{the same released GRAM checkpoint} in
three environments. The spread is $2.3$ R@1 between the published number and
ours, which is an order of magnitude larger than our evaluation-noise floor and
comparable to the margins the field reports as contributions. This is why the
main paper compares only artifacts evaluated in one environment and retains
published numbers for reference only.

The offset is also not sign-consistent, so it cannot be corrected by a constant.
HyperGRAM trained from its authors' unmodified code at their recipe measures
$-2.6$, $-2.4$ and $-4.5$ R@1 against its published numbers on MSR-VTT, DiDeMo
and ActivityNet, but $+9.9$ on VATEX; GRAM's released checkpoint measures
$-2.3$ on MSR-VTT and $+6.5$ on VATEX.

\begin{table}[t]
\centering
\caption{Evaluation-environment audit. The first three rows evaluate the same
released GRAM checkpoint in three environments; the fourth is our own retraining
of GRAM evaluated in ours. MSR-VTT T-VAS, two-stage R@1.}
\label{tab:envaudit}
\small
\begin{tabular}{llr}
\toprule
Weights & Evaluation environment & R@1 \\
\midrule
GRAM, released checkpoint & as published & 54.8 \\
GRAM, released checkpoint & HyperAlign & 53.4 \\
GRAM, released checkpoint & ours & 52.5 \\
GRAM, our retraining & ours & 52.4\,\tiny{$\pm$0.21} \\
\bottomrule
\end{tabular}
\end{table}

\subsection{Statistical Support}
\label{sec:stats}

Every count stated in the main paper is accompanied by an exact test here.

\paragraph{Missing-modality sweep.}
\ours{} exceeds GRAM in 25 of 25 rate$\times$benchmark cells; exact one-sided
sign test $p = 2.98\times10^{-8}$. Cells within a benchmark share nested masks
and are therefore not independent, so we also report the conservative version
that treats each benchmark as a single unit: 5 of 5, $p = 0.031$.

\paragraph{Aggregation gain.}
12 of the 14 measurable baseline method$\times$benchmark cells are negative;
against a null in which fusion helps or hurts with equal probability, exact
one-sided $p = 0.0065$. Method rows share checkpoints across benchmarks, so
these cells are not fully independent either; the count itself is the primary
evidence and the test is reported for calibration.

\paragraph{Seed variance versus baseline gaps.}
Table~\ref{tab:seedci} gives the three-seed mean, standard deviation and 95\%
interval for \ours{} against each single-run baseline. Released checkpoints are
$n{=}1$ by nature, so no two-sample test exists; ``outside'' states only that
the observed ordering is not explained by \ours{}'s seed variance, and says
nothing about the baseline's own variance.

\begin{table}[t]
\centering
\caption{\ours{} seed statistics against single-run baselines (two-stage
T$\rightarrow$V R@1). Baselines are released checkpoints or single training
runs, so only \ours{} has a variance estimate.}
\label{tab:seedci}
\small
\setlength{\tabcolsep}{3pt}
\begin{tabular}{lccccc}
\toprule
Benchmark & \ours{} $\pm$ sd & 95\% CI & GRAM$^{\star}$ & PMRL$^{\star}$ & HyperGRAM$^{\dagger}$ \\
\midrule
MSR-VTT & 54.6 $\pm$ 0.2 & [54.1, 55.2] & 52.5 & 54.3 & 54.0 \\
DiDeMo & 51.2 $\pm$ 0.3 & [50.6, 51.9] & 50.7 & 52.5 & 48.9 \\
ActivityNet & 55.9 $\pm$ 0.2 & [55.6, 56.3] & 56.3 & 54.1 & 53.7 \\
VATEX & 90.6 $\pm$ 0.1 & [90.3, 90.9] & 90.0 & 89.6 & 89.8 \\
AudioCaps & 35.0 $\pm$ 0.4 & [34.0, 36.0] & 32.2 & 34.4 & -- \\
\bottomrule
\end{tabular}
\end{table}

\subsection{Trainable-Parameter Accounting}
\label{sec:params}

We report a trainable-parameter count only for configurations we trained
ourselves, and describe the baselines as pretraining the full backbone rather
than assigning them a number. Two measurements motivate this.

First, the two available counting rules are not interchangeable. For a low-rank
run the honest count is the set of parameters the optimizer updated; for a
released checkpoint the only available count is every floating-point tensor in
the file, which also includes buffers and heads that never receive gradients.
Our own full-fine-tuning control is measurable both ways and reads
$1{,}242{,}890{,}226$ by the first rule against $1{,}397{,}367{,}145$ by the
second---an $11\%$ gap on one identical model. A count from one rule cannot sit
in a column beside a count from the other.

Second, the checkpoint rule also counts trunk artifacts that are not part of the
method. A key-level comparison of the checkpoints shows that our HyperGRAM
training arm and our full-fine-tuning arm each carry a $720{,}896$-parameter
depth projection head that the shared trunk builds unconditionally and that is
never invoked in the modality configurations evaluated here; the released GRAM
checkpoint does not contain it. Reporting that as HyperGRAM's parameter count
would attribute our trunk's unused weights to their method.

The comparison the paper makes is therefore stated in the form it can support:
\ours{} trains 4{,}763{,}650 parameters, and every baseline pretrains the full
backbone.

% =====================================================================
\section{Limitations}
\label{sec:limits}
% =====================================================================

\paragraph{One retrieval direction.}
The aggregation is conditioned on the text query, so the mechanism has no
purchase on video-to-text retrieval, where \ours{} is mid-field. We report that
direction rather than omit it.

\paragraph{Short queries.}
Query conditioning requires a query with enough content to condition on. On
VGGSound, where queries are single class names, the uniform-weights control is
stronger than the query-weighted model (Section~\ref{sec:vggsound}).

\paragraph{One backbone.}
All results use the VAST backbone, because that is what makes the comparison to
GRAM, PMRL and HyperGRAM an artifact-level comparison rather than a comparison
across training pipelines. Whether the aggregation-gain deficit reproduces on
other multimodal backbones is untested here.

\paragraph{Two-stage ceiling.}
Our representation-level gains are compressed by the reranking stage, and under
masking they are erased by it (Section~\ref{sec:twostage}). We diagnose this
bottleneck but do not fix it; a reranker trained on incomplete modality sets is
the obvious next step and is out of scope here.

\paragraph{Corpus scale.}
Following the compared methods, continued pretraining uses a 150k-clip subset
rather than full VAST-27M, so all conclusions are drawn in the small-budget
adaptation regime that this family of methods operates in.

\end{document}
